\subsection{Short integral generators from the actual theta series}
\label{sec:actual-short-theta-generators}

The spanning result becomes quantitative without bounding its spanning
coefficients.  We choose short integral harmonic polynomials, estimate
their actual raw theta series, and multiply each series by a small
element of its prescribed determinant ideal.  The last multiplication
cancels the lattice index in the analytic estimate.  All the fields,
levels, lattices and theta normalizations remain those of the construction
\cite{theta-short-source}.

Use the actual sized good cover and all the original
\texttt{CovGmin.Hyp} assumptions.  In particular put
\[
 f=[E:\mathbb Q]/2,\quad r=p+q,\quad K=c.F,\quad
 X=frL_b,\quad d=\mathrm{DF}(\mathrm{Dw},a),\quad
 \kappa=d+\frac n2fr.
\]
Thus $p,q\ge1$, $f>2$, $n$ is even, $n\ge8fr+4$, and the
original positive-weight condition holds.  Retain
\[
 \begin{gathered}
 L_b\ge1,\quad \log M\le30L_b,\quad
 \operatorname{fieldSz}(E)\le3L_b,\\
 \log(fr)\le L_b,\quad \log\kappa\le20L_b,\quad
 \operatorname{fieldSz}(K)\le o\le391X^2.
 \end{gathered}
\]
The datum, principal-level laws, common component matrix and form
constant, full trivial arithmetic character, and normalized dictionary
are the actual ones constructed earlier.  The condition on $K$ retains
both its degree and root discriminant; $K$ is not fixed as the table
varies.

\begin{theorem}[Short integral theta spans]
\label{thm:actual-short-theta-generators}
The original cusp space $Y.V(\mathrm{wtZ}(a,n))$ has a finite
$K$-spanning family in its original integral cusp lattice whose global
Petersson norm and every original component Petersson norm, at every
embedding of $K$, are at most
\begin{equation}\label{eq:actual-short-theta-bound}
                         \exp(30000X^2\kappa).
\end{equation}
Every member has the form $Zs_\Lambda\theta_{\Lambda,N,x_0,P}$,
where $\theta$ is the actual raw theta section, $P$ belongs to the
original $\mathrm{harmL}$, $N$ has the original odd level condition,
$s_\Lambda\ne0$ satisfies the original determinant-ideal inclusion,
and
\[
             Z=M^{frn+d}|\operatorname{disc}E|^d
\]
is the same clearing integer for the whole family.  No claim that this
family generates the integral module is required or made.
\end{theorem}

\begin{proof}
\emph{Short vectors in a unimodular Hermitian lattice.}
Let $\Lambda\subset E^n$ be one of the positive definite unimodular
lattices used in the actual theta span.  It has $2fn$ rationally
independent vectors $v_j\in\Lambda$ satisfying, for every embedding,
\begin{equation}\label{eq:actual-short-lattice-vectors}
             \|\tau v_j\|\le B_0,
             \qquad B_0=2f^{3/2}n^2\operatorname{rd}(E).
\end{equation}
Here is a proof that also retains the different.  For the Euclidean
form $B(u,v)=\operatorname{Tr}_{E/\mathbb Q}h(u,v)$, an element of
the trace-dual lattice has $h(u,\Lambda)\subset\mathfrak D_E^{-1}$.
Hermitian self-duality gives $\mathfrak D_Eu\subset\Lambda$, hence
$h(u,u)\in\mathfrak D_E^{-2}$.  If $u\ne0$, the ideal norm and
arithmetic--geometric mean give
\[
 B(u,u)\ge 2f\,|N_{E/\mathbb Q}h(u,u)|^{1/(2f)}
              \ge\frac{2f}{\operatorname{rd}(E)^2}.
\]
These are trace-dual ideals, including at ramified primes.

We use the elementary Euclidean transference bound
$\lambda_m(L)\lambda_1(L^*)\le m^2$.  For completeness,
Minkowski applied to cubes gives
$\lambda_1(L)\lambda_1(L^*)\le m$.  Choose a primitive short
dual vector $y$, of length at most $m/\lambda_1(L)$, and apply
induction to $L\cap y^\perp$.  Its minimum is at least
$\lambda_1(L)$.  Extend each of the resulting dual vectors to
$L^*$ by adding a multiple of $y$, with coefficient of absolute
value at most $1/2$ after rounding.  The squared length is at most
$((m-1)^4+m^2/4)/\lambda_1(L)^2\le m^4/\lambda_1(L)^2$.
Together with $y$ these vectors are independent.  Applying the result
to the trace dual with $m=2fn$, and using
$2\|\tau v\|^2\le B(v,v)$, proves
\eqref{eq:actual-short-lattice-vectors}.  This argument requires no
free $\mathcal O_E$-basis of $\Lambda$.

\emph{Short integral harmonic polynomials.}
The actual harmonic highest-weight space of total degree $d$ is
spanned over $K$ by integral polynomials $P$ with
\begin{equation}\label{eq:actual-short-harmonic-bound}
 \log\sum_\alpha|\sigma(P_\alpha)|\le140rL_*d
 \quad(\sigma:K\hookrightarrow\mathbb C),
 \qquad
 L_*=2+3L_b+\log(1+d)+\log n\le50L_b.
\end{equation}
We give the finite arithmetic construction.  Put $b=(2r+1)d$ and
choose three arrays of vectors $U,A,W$ from the grid
$\{\sum_j a_jv_j:0\le a_j\le b,\ a_j\in\mathbb Z\}$, of
lengths $p,p,q$.  Every grid vector has embedding norm at most
$R=(2fn)bB_0$.  With $H_{ab}=h(A_b,U_a)$ define
\[
 V_c=(\det H)W_c-
       \sum_b A_b\bigl(\operatorname{adj}(H)
                         (h(W_c,U_a))_a\bigr)_b.
\]
Here the original Hermitian convention is explicitly
$h(x,y)=\sum_kx_k\overline{y_k}$, linear in the first variable
(\texttt{EmbeddingHermSelf}, lines 14--26).  Thus the coefficients
in this formula are not conjugated.
Then $V_c\in\Lambda$ and $h(V_c,U_a)=0$, including when $H$ is
singular.  The integral leading-minor products made from $U$ and $V$
are the original harmonic row-highest polynomials, after the fixed
letter dictionary.  Indeed contractions between the two sets of linear
forms are zero, and their coefficients belong to the original symmetric
integral algebra of $\Lambda$.

These products span the full harmonic space.  One elementary
highest-weight proof is as follows.  Write the two arrays of variables
as $z_{ik}$ and $y_{jk}$, with $p$ and $q$ rows.  For row partitions
$\alpha,\beta$ and $p+q\le n$, put
\[
 \nu=(\alpha_1,\ldots,\alpha_p,0,\ldots,0,
                         -\beta_q,\ldots,-\beta_1).
\]
Row-highest invariance, in addition to the row degrees, implies that
the degree of any monomial in a set $S$ of $\ell$ columns of the first
array is at most $\sum_{i\le\min(\ell,p)}\alpha_i$.
For $\ell<p$, prove this on the open set where the first $\ell$ rows
restricted to $S$ are independent.  Adding earlier rows to later rows
eliminates the entries of $S$ in every later row and leaves the
polynomial unchanged.  Scaling the columns $S$ by $t$ leaves these
elimination coefficients unchanged, so the resulting $t$-degree is
at most $\alpha_1+\cdots+\alpha_\ell$.  The identity on this dense
open set proves the bound everywhere.  For $\ell\ge p$ it is the
total degree bound.  The same argument applies to $y$ and $\beta$.
Consequently every column weight $\mu$ satisfies
\[
 \sum_{k\le\ell}\mu_k
 \le\sum_{i\le\min(\ell,p)}\alpha_i-|\beta|
       +\sum_{j\le\min(n-\ell,q)}\beta_j
 =\sum_{k\le\ell}\nu_k.
\]
Thus every dominant column-highest weight is dominated by $\nu$.

For clarity, the exact differential-operator identity used next is
\[
 \Omega_{\rm col}=\Omega_z+(n-p)E_z
                   +\Omega_y+(n-q)E_y
                     -2\sum_{i,j}R_{ij}\Delta_{ij},
\]
where $E_z,E_y$ are total degrees,
$R_{ij}=\sum_k z_{ik}y_{jk}$,
$\Delta_{ij}=\sum_k\partial_{z_{ik}}\partial_{y_{jk}}$,
the row Casimirs are $\sum_{i,j}(\sum_kz_{ik}\partial_{z_{jk}})
(\sum_kz_{jk}\partial_{z_{ik}})$ and its $y$ analogue, and
the column Casimir is $\sum_{k,l}L_{kl}L_{lk}$ with
$L_{kl}=\sum_i z_{ik}\partial_{z_{il}}-
\sum_j y_{jl}\partial_{y_{jk}}$.
Expansion and the product rule prove the identity.
The row Casimir on a highest vector of weight $\alpha$ is
$\sum_i\alpha_i(\alpha_i+p+1-2i)$, by commuting each raising
operator past its lowering operator; similarly for $\beta$.
Harmonicity therefore makes the column Casimir equal to
$\sum_i\nu_i(\nu_i+n+1-2i)$.  This value is strictly increasing
under strict dominance: summation by parts has positive coefficients
$(\nu_j-\nu_{j+1})+(\mu_j-\mu_{j+1})+2$.
Consequently $\mu=\nu$.  The monomial degrees then exclude all
columns except the first $p$ in the first block and the last $q$ in
the second.  On the open set where the successive leading minors
are nonzero, row and column unipotent elimination reduces these two
square blocks to diagonal matrices.  The diagonal weights fix the
single product of leading minors, with exponents
$\alpha_i-\alpha_{i+1}$ and $\beta_j-\beta_{j+1}$.
Polynomial identity extends this uniqueness off the open set.
In the Fock inner product, column raising
and lowering operators are adjoints.  The orthogonal complement of
the span of this product's column orbit, if nonzero, would contain
another highest vector and hence this same product, a contradiction.
This proves orbit spanning.  The assertion descends from $\mathbb C$
to $K$ because all the defining equations and operators are linear
over $K$.

The adjugate construction parametrizes this orbit on the open set of
independent $U$: take $A$ dual to $U$, so that $H=1$, and take any
prescribed orthogonal $V=W$.  The distinct CM embeddings and their
conjugates are independent after scalar extension; the rational basis
$(v_j)$ therefore supplies all these parameters simultaneously.
As a polynomial in the grid coordinates each minor product has degree
at most $(2p+1)d\le b$.  Interpolation on $\{0,\ldots,b\}$ in
each coordinate proves that the stated integral grid already spans.
Its cardinality is at most
$((2r+1)d+1)^{2fn(2p+q)}$.

The adjugate formula gives
$\|\tau V_c\|\le(p+1)!R^{2p+1}$.
Hadamard expansion and the coefficient $\ell^1$ norm bound every
minor product by
\[
 \left(r\sqrt n\,(r+1)!R^{2r+1}\right)^d.
\]
Using $r\le n$, $\log\operatorname{rd}(E),\log f,
\log(1+d)\le2L_*$ and $\log n\le L_*$ gives
\eqref{eq:actual-short-harmonic-bound}.  If $d=0$, use $\{1\}$
when the harmonic space is nonzero, and the empty family otherwise.
Thus no logarithm of a zero degree is taken.  For the source's global
coefficient norm,
\[
 \operatorname{polyNorm}(P)
 \le\sqrt{[K:\mathbb Q]}\max_\sigma\sum_\alpha|\sigma P_\alpha|
 \le\exp(20000X^2\kappa).
\]
The term $\tfrac12\log[K:\mathbb Q]\le o/2$ is retained in this
inequality.

\emph{The norm of the actual raw theta section.}
Write $m_\Lambda=[\Lambda:\Lambda\cap\mathcal O_E^n]$.  For a
homogeneous original $P$ with $\operatorname{polyNorm}(P)\le\rho$,
the actual section has normalized point norm bounded by
\begin{equation}\label{eq:actual-raw-theta-uniform-norm}
 \operatorname{homN}_\sigma(\theta;z)
       \le m_\Lambda^{r/2}\rho\exp(9000X^2\kappa).
\end{equation}
We verify the decisive analytic estimates and the index.

First, a nonzero bounded cusp form has a frame at which its point norm
is at least half its supremum.  Normalize that frame orthogonally.
At this frame the original period lattice has minimum at least
\begin{equation}\label{eq:actual-halfmax-minimum}
                                \lambda=\exp(-40L_b).
\end{equation}
Here no uniform fundamental domain is assumed.  The actual principal
level transvections give a direct proof.  On the integral Hermitian
lattice $L'$ with $\Lambda_{\rm per}=\mathrm{AA}L'$, the majorant
of a vector smaller than one at every embedding forces its Hermitian
norm to vanish, by integrality and the product formula.  For such an
isotropic vector use its rank-one unipotent flow with complex parameter
$\zeta$.  If all majorants are at most $\epsilon$, the orbit remains
in the domain for $\Im\zeta>-T$, $T=(k_0\epsilon)^{-1}$.
The original arithmetic group makes the evaluated form periodic of
period $M_*$, with
\[
             k_0=2\xi_R,\qquad M_*k_0=2M\sqrt{d_H},
\]
where $d_H$ is the Host's original quadratic parameter.  The trivial
arithmetic character gives genuine periodicity.  Cuspidality gives
vanishing at $\Im\zeta=+\infty$.  On $\Im\zeta=-T/2$ the
frame-norm distortion is at most $2^{w/2}$, where
$w=3\kappa$ bounds the sum of the weight spreads.  The periodic
Schwarz lemma and the half-supremum choice would imply
\[
                  1\le2\,2^{w/2}\exp(-\pi T/M_*).
\]
This is impossible for
\[
 \epsilon\le\epsilon_{\rm lev}
 =\min\left\{1,
       \frac{\pi}{2M\sqrt{d_H}(fw+2)\log2}\right\}.
\]
The actual trace-dual identity
$(\xi_0)\mathrm{AA}\,\overline{\mathrm{AA}}\mathfrak D_E=1$
and Minkowski give $b\ne0$ with $b\mathrm{AA}\subset\mathcal O_E$
and $|\tau b|^2\le\xi_R\operatorname{rd}(E)^2$.  Applying the
majorant conclusion to $b\Lambda_{\rm per}\subset L'$ gives
$\lambda_1(\Lambda_{\rm per})^2
 \ge2\epsilon_{\rm lev}/\operatorname{rd}(E)^2$.
Finally $\log\sqrt{d_H}\le\log\operatorname{rd}(E)\le3L_b$
and the displayed size assumptions give
$\epsilon_{\rm lev}\ge e^{-57L_b}$ and hence
\eqref{eq:actual-halfmax-minimum}.

For completeness the remaining theta bound is a Gaussian estimate.
The tensor period lattice $\Lambda\otimes\Lambda_{\rm per}$ has
minimum at least $\lambda/B_0$: contraction with one of the vectors
in \eqref{eq:actual-short-lattice-vectors} is a nonzero original
period and has norm at most $B_0$ times the tensor's norm.
On a principal torus of complex dimension $g=frn$ and period minimum
$\lambda'$, Gaussian submean on a ball of radius $\sqrt{2g/\pi}$
and Euclidean packing give
\[
 \|s\|_\infty\le\sqrt2
   \left(1+\frac{4\sqrt{2g}}{\sqrt\pi\lambda'}\right)^g
   \|s\|_2.
\]
The Gaussian mass of this ball is at least one half; its translated
lattice multiplicity is at most
$(1+4\sqrt{2g}/(\sqrt\pi\lambda'))^{2g}$.
This proves the inequality with the original average $L^2$ norm.
The finite isogeny sum defining the projected vacuum has squared
$L^2$ norm equal to the number of its cosets times the old squared
norm.  This is finite Heisenberg character orthogonality on their
common cover.  Dividing by the original square root of that coset
number gives norm one, without any assumption on the relative index.

Cauchy's estimate for ordinary derivatives of total degree $d$, with
the radius $\sqrt{\alpha_j/\pi}$ in coordinate $j$, gives the
factor $(\pi e d)^{d/2}$.  The actual derivative directions contribute
$t^{d/2}$, $t=2\xi_R$; no factorial is inserted into the harmonic
operator.  Consequently at the chosen frame the actual raw value is
at most
\[
 m_\Lambda^{r/2}\Bigl(\sum_\alpha|\sigma P_\alpha|\Bigr)
 (\pi e t d)^{d/2}\sqrt2
 \left(1+\frac{4\sqrt{2frn}\,B_0}{\sqrt\pi\lambda}\right)^{frn}.
\]
The normalized finite Fourier mean has coefficients of modulus one;
its full cardinality cancels its denominator.  Combining the actual
raw prefactor with the normalized vacuum leaves exactly
$m_\Lambda^{r/2}$, as in the raw-span calculation.  The original
Lie normalization is $(\pi t)^{-d/2}$ and cancels its matching
derivative factor.  For $d=0$ these factors are one.
Now $d\le\kappa$, $frn\le2\kappa$,
$\#\operatorname{supp}P\le(d+1)^{frn}$, and
$\sum|\sigma P_\alpha|\le\sqrt{\#\operatorname{supp}P}\,\rho$.
The logarithms of these factors, including the half-supremum factor,
are bounded by $9000X^2\kappa$.  For example
$\log n\le21L_b$ and the logarithm of the packing parenthesis
is at most $101L_b$; their sum is already less than
$300\kappa L_b\le9000X^2\kappa$.
This proves \eqref{eq:actual-raw-theta-uniform-norm} on every component
and at every embedding with the original normalization.

\emph{Cancelling the index while preserving integrality.}
Let $\mathfrak a=\det(\Lambda\cap\mathcal O_E^n)(\det\Lambda)^{-1}$.
It is integral and $N\mathfrak a=m_\Lambda$.  In $K$ put
\[
 \mathfrak B=\prod_{\tau:E\hookrightarrow\mathbb C}
                 (\tau\mathfrak a)^{\mathrm{Dw.mult}(\tau)}.
\]
Because $\sum_\tau\mathrm{Dw.mult}(\tau)=fr$, its norm is
$N_K\mathfrak B=m_\Lambda^{r[K:\mathbb Q]/2}$.
Minkowski in the inverse fractional ideal, using a product of equal
complex discs, gives a nonzero $s_\Lambda\in\mathfrak B^{-1}$ with
\begin{equation}\label{eq:actual-small-determinant-scalar}
        |\sigma s_\Lambda|\,m_\Lambda^{r/2}
                  \le e^{\operatorname{fieldSz}(K)}
                  \qquad\hbox{for every }\sigma.
\end{equation}
Indeed, writing $[K:\mathbb Q]=2h$, the discs of radius
$e^{\operatorname{fieldSz}(K)}m_\Lambda^{-r/2}$ have volume
$\pi^h e^{2h\operatorname{fieldSz}(K)}m_\Lambda^{-rh}$,
whereas the required Minkowski threshold is
$2^h\sqrt{|\operatorname{disc}K|}m_\Lambda^{-rh}$.
The former is strictly larger since
$\operatorname{fieldSz}(K)\ge\log\operatorname{rd}(K)$.
Membership in $\mathfrak B^{-1}$ is precisely the original
determinant-ideal inclusion.  Thus this is the same scalar in both
the integral and analytic arguments.

The established actual denominator theorem\cite{current-construction-record} makes
$Zs_\Lambda\theta$ integral, with $\log Z\le67X^2\kappa$.
Equations \eqref{eq:actual-short-harmonic-bound},
\eqref{eq:actual-raw-theta-uniform-norm} and
\eqref{eq:actual-small-determinant-scalar} give total logarithmic
bound
\[
        (20000+9000+391+67)X^2\kappa<30000X^2\kappa.
\]
The original global and component Petersson measures have mass at
most one, so the same bound holds for their $L^2$ norms.  The actual
raw-span theorem and the harmonic grid span show that these series
span the original cusp space.  Multiplication of individual generators
by nonzero $Zs_\Lambda$ preserves that $K$-span.  Finite dimensionality
then permits selection of a finite spanning subset without increasing
any norm.  This proves the theorem.
\end{proof}

\paragraph{The precise downstream gain.}
This supplies the original \texttt{ExistsGoodCoverAtMinimalLevel}
conclusion on the actual cover.  Its argument does not use the
\texttt{PublishedHodgeBundleTheorems} package appearing in the source
wrapper.  The next weight-shift supplier must still be applied to the
same cover and its actual integral flat-line and key-unit twists.
Once its integral $K$-spanning set has the displayed component norm,
the $h_0,h_{\rm sh}$ consumers in
\texttt{ExistsDegreeParametersBound} require only a finite subset.
The intervening conversion to generators of the whole integral module,
with its additional rank and Tamagawa inputs, is unnecessary for these
consumers.  This observation does not assert that the twists or their
flat degree have been supplied by the present theorem.
